\begin{resumo}[Abstract]
 \begin{otherlanguage*}{english}
   The contemporary university environment has seen a significant increase in cases related to health problems among students, including mental health issues, sedentary lifestyles, academic stress, and habits that are detrimental to well-being. These factors directly impact student performance and highlight the need for effective monitoring and intervention mechanisms. Therefore, this work aims to develop a mobile application aimed at collecting and analyzing data on mood, habits, and general health within higher education institutions (HEIs). The proposal aims to help guide institutional health promotion policies and provide students with ways to seek help within the HEI. The methodology adopted involves a literature and document review on the use of health technologies, analysis of similar applications, and the administration of questionnaires to the academic community to identify the application's functional and non-functional requirements. Therefore, the next phase will begin software development, including the implementation of features, to develop a means to assist students in adopting healthy habits and promoting mental health within the academic environment.

   \vspace{\onelineskip}
 
   \noindent 
   \textbf{Key-words}: Mobile application, Healthy Habits, Mental Health.
 \end{otherlanguage*}
\end{resumo}
